\documentclass[aspectratio=169]{ctexbeamer}

\usepackage{tikz}
\usetikzlibrary{positioning}
\usepackage{booktabs}
\usepackage{fontawesome5}

%\usepackage[UTF8]{ctex}
\usepackage{listings}
\usepackage{xcolor}
\usepackage{hyperref}

\usetheme{Madrid}
\usecolortheme{default}

\definecolor{codebg}{RGB}{245,245,245}
\definecolor{codekw}{RGB}{0,92,175}
\definecolor{codestr}{RGB}{163,21,21}
\definecolor{codecomment}{RGB}{0,128,0}

\lstset{
  backgroundcolor=\color{codebg},
  basicstyle=\ttfamily\small,
  keywordstyle=\color{codekw}\bfseries,
  stringstyle=\color{codestr},
  commentstyle=\color{codecomment},
  numbers=left,
  numberstyle=\tiny\color{gray},
  frame=single,
  breaklines=true,
  columns=fullflexible,
  keepspaces=true,
  showstringspaces=false,
  tabsize=4
}

\title{Make 入门讲座}
\subtitle{用 Makefile 自动化构建与任务管理}
\author{Your Name}
\date{\today}

\begin{document}

%------------------------------------------------
\begin{frame}
  \titlepage
\end{frame}

%------------------------------------------------
\begin{frame}{本讲内容}
  \tableofcontents
\end{frame}

%------------------------------------------------
\section{为什么需要 Make}

\begin{frame}[fragile]{问题：手动编译很麻烦}
假设有一个简单的 C 项目：

\begin{itemize}
  \item \texttt{main.c}
  \item \texttt{hello.c}
  \item \texttt{hello.h}
\end{itemize}

手动编译可能需要输入：

\begin{lstlisting}[language=bash]
gcc -c main.c
gcc -c hello.c
gcc -o app main.o hello.o
\end{lstlisting}

问题：

\begin{itemize}
  \item 文件多了以后命令很长
  \item 修改一个文件后不知道该重新编译哪些文件
  \item 容易输错命令
  \item 不利于团队协作
\end{itemize}
\end{frame}

%------------------------------------------------
\begin{frame}{Make 能做什么？}
\textbf{Make} 是一种自动化构建工具。

它可以根据文件之间的依赖关系，自动判断哪些任务需要重新执行。

\begin{block}{核心思想}
如果目标文件不存在，或者依赖文件比目标文件更新，则重新执行命令。
\end{block}

常见用途：

\begin{itemize}
  \item 编译 C / C++ 项目
  \item 自动运行测试
  \item 清理临时文件
  \item 生成文档
  \item 执行脚本任务
\end{itemize}
\end{frame}

%------------------------------------------------
\section{Makefile 基础}

\begin{frame}[fragile]{Makefile 的基本结构}
Make 使用一个名为 \texttt{Makefile} 的文件描述构建规则。

基本格式：

\begin{lstlisting}[language=make]
target: dependencies
	command
\end{lstlisting}

含义：

\begin{itemize}
  \item \texttt{target}：目标，要生成的文件或任务名
  \item \texttt{dependencies}：依赖，生成目标需要的文件
  \item \texttt{command}：命令，用来生成目标
\end{itemize}

\begin{alertblock}{注意}
命令行前面必须是 \textbf{Tab}，不能是普通空格。
\end{alertblock}
\end{frame}

%------------------------------------------------
\begin{frame}[fragile]{一个最简单的 Makefile}
创建文件 \texttt{Makefile}：

\begin{lstlisting}[language=make]
hello:
	echo "Hello, Make!"
\end{lstlisting}

在终端执行：

\begin{lstlisting}[language=bash]
make
\end{lstlisting}

输出类似：

\begin{lstlisting}[language=bash]
echo "Hello, Make!"
Hello, Make!
\end{lstlisting}

说明：

\begin{itemize}
  \item 默认执行 Makefile 中的第一个目标
  \item 这里第一个目标是 \texttt{hello}
\end{itemize}
\end{frame}

%------------------------------------------------
\begin{frame}[fragile,shrink]{指定目标执行}
如果 Makefile 中有多个目标：

\begin{lstlisting}[language=make]
hello:
	echo "Hello"

bye:
	echo "Bye"
\end{lstlisting}

可以指定目标：

\begin{lstlisting}[language=bash]
make hello
make bye
\end{lstlisting}

如果只输入：

\begin{lstlisting}[language=bash]
make
\end{lstlisting}

则执行第一个目标：

\begin{lstlisting}[language=make]
hello
\end{lstlisting}
\end{frame}

%------------------------------------------------
\section{简单 C 项目示例}

\begin{frame}[fragile]{示例项目结构}
项目结构如下：

\begin{lstlisting}
project/
|-- Makefile
|-- main.c
|-- hello.c
`-- hello.h
\end{lstlisting}

文件 \texttt{hello.h}：

\begin{lstlisting}[language=C]
#ifndef HELLO_H
#define HELLO_H

void say_hello(void);

#endif
\end{lstlisting}
\end{frame}

%------------------------------------------------
\begin{frame}[fragile]{示例代码：hello.c}
文件 \texttt{hello.c}：

\begin{lstlisting}[language=C]
#include <stdio.h>
#include "hello.h"

void say_hello(void) {
    printf("Hello, Make!\n");
}
\end{lstlisting}
\end{frame}

%------------------------------------------------
\begin{frame}[fragile]{示例代码：main.c}
文件 \texttt{main.c}：

\begin{lstlisting}[language=C]
#include "hello.h"

int main(void) {
    say_hello();
    return 0;
}
\end{lstlisting}
\end{frame}

%------------------------------------------------
\begin{frame}[fragile]{手动编译与运行}
手动编译：

\begin{lstlisting}[language=bash]
gcc -c main.c
gcc -c hello.c
gcc -o app main.o hello.o
./app
\end{lstlisting}

运行结果：

\begin{lstlisting}[language=bash]
Hello, Make!
\end{lstlisting}
\end{frame}

%------------------------------------------------
\begin{frame}[fragile,shrink]{第一版 Makefile}
\begin{lstlisting}[language=make]
app: main.o hello.o
	gcc -o app main.o hello.o

main.o: main.c hello.h
	gcc -c main.c

hello.o: hello.c hello.h
	gcc -c hello.c
\end{lstlisting}

执行：

\begin{lstlisting}[language=bash]
make
\end{lstlisting}

Make 会分析依赖关系：

\begin{itemize}
  \item \texttt{app} 依赖 \texttt{main.o} 和 \texttt{hello.o}
  \item \texttt{main.o} 依赖 \texttt{main.c} 和 \texttt{hello.h}
  \item \texttt{hello.o} 依赖 \texttt{hello.c} 和 \texttt{hello.h}
\end{itemize}
\end{frame}

%------------------------------------------------
\begin{frame}[fragile,shrink]{Make 的增量构建}
第一次执行：

\begin{lstlisting}[language=bash]
make
\end{lstlisting}

会编译所有文件。

如果再次执行：

\begin{lstlisting}[language=bash]
make
\end{lstlisting}

可能看到：

\begin{lstlisting}[language=bash]
make: 'app' is up to date.
\end{lstlisting}

含义：

\begin{itemize}
  \item 目标文件已经是最新的
  \item 不需要重复编译
\end{itemize}

如果修改 \texttt{hello.c}，再执行 \texttt{make}：

\begin{itemize}
  \item 只会重新编译 \texttt{hello.o}
  \item 然后重新链接生成 \texttt{app}
\end{itemize}
\end{frame}

%------------------------------------------------
\section{常用写法}

\begin{frame}[fragile,shrink]{使用变量}
Makefile 可以定义变量，减少重复代码。

\begin{lstlisting}[language=make]
CC = gcc
CFLAGS = -Wall -O2
TARGET = app
OBJS = main.o hello.o

$(TARGET): $(OBJS)
	$(CC) -o $(TARGET) $(OBJS)

main.o: main.c hello.h
	$(CC) $(CFLAGS) -c main.c

hello.o: hello.c hello.h
	$(CC) $(CFLAGS) -c hello.c
\end{lstlisting}

优点：

\begin{itemize}
  \item 修改编译器更方便
  \item 编译参数统一管理
  \item Makefile 更易维护
\end{itemize}
\end{frame}

%------------------------------------------------
\begin{frame}[fragile]{伪目标：.PHONY}
有些目标不是文件，而是一个任务，例如 \texttt{clean}。

\begin{lstlisting}[language=make]
.PHONY: clean

clean:
	rm -f app main.o hello.o
\end{lstlisting}

执行：

\begin{lstlisting}[language=bash]
make clean
\end{lstlisting}

作用：

\begin{itemize}
  \item 删除编译生成的文件
  \item 方便重新构建
\end{itemize}

为什么需要 \texttt{.PHONY}？

\begin{itemize}
  \item 告诉 Make：\texttt{clean} 不是一个真实文件
  \item 即使目录中有名为 \texttt{clean} 的文件，也照样执行该任务
\end{itemize}
\end{frame}

%------------------------------------------------

\begin{frame}[fragile,shrink]{自动变量}
Make 提供了一些自动变量，可以让规则更简洁。

\begin{center}
\begin{tabular}{c|l}
\textbf{变量} & \textbf{含义} \\
\hline
\texttt{\$@} & 当前目标文件 \\
\texttt{\$<} & 第一个依赖文件 \\
\texttt{\$\textasciicircum} & 所有依赖文件 \\
\end{tabular}
\end{center}

示例：

\begin{lstlisting}[language=make]
app: main.o hello.o
	gcc -o $@ $^

main.o: main.c hello.h
	gcc -c $<

hello.o: hello.c hello.h
	gcc -c $<
\end{lstlisting}

解释：

\begin{itemize}
  \item \texttt{\$@} 在第一条规则中表示 \texttt{app}
  \item \texttt{\$\textasciicircum\{\}} 表示 \texttt{main.o hello.o}
  \item \texttt{\$<} 表示第一个依赖文件
\end{itemize}
\end{frame}

%------------------------------------------------
\begin{frame}[fragile]{模式规则}
如果很多 \texttt{.c} 文件都要编译成 \texttt{.o} 文件，可以写成模式规则。

\begin{lstlisting}[language=make]
CC = gcc
CFLAGS = -Wall -O2
TARGET = app
OBJS = main.o hello.o

$(TARGET): $(OBJS)
	$(CC) -o $@ $^

%.o: %.c
	$(CC) $(CFLAGS) -c $< -o $@
\end{lstlisting}

含义：

\begin{itemize}
  \item \texttt{\%.o: \%.c} 表示任意 \texttt{.o} 文件都可以由对应的 \texttt{.c} 文件生成
  \item 例如：\texttt{main.o} 由 \texttt{main.c} 生成
  \item 例如：\texttt{hello.o} 由 \texttt{hello.c} 生成
\end{itemize}
\end{frame}

%------------------------------------------------
\begin{frame}[fragile,shrink]{加入 clean 和 run}
一个更实用的 Makefile：

\begin{lstlisting}[language=make]
CC = gcc
CFLAGS = -Wall -O2
TARGET = app
OBJS = main.o hello.o

.PHONY: all clean run

all: $(TARGET)

$(TARGET): $(OBJS)
	$(CC) -o $@ $^

%.o: %.c
	$(CC) $(CFLAGS) -c $< -o $@

run: $(TARGET)
	./$(TARGET)

clean:
	rm -f $(TARGET) $(OBJS)
\end{lstlisting}

常用操作：

\begin{lstlisting}[language=bash]
make
make run
make clean
\end{lstlisting}
\end{frame}

%------------------------------------------------
\section{常用命令与调试}

\begin{frame}[fragile]{常用 make 命令}
\begin{lstlisting}[language=bash]
make
\end{lstlisting}
执行默认目标。

\begin{lstlisting}[language=bash]
make clean
\end{lstlisting}
执行清理任务。

\begin{lstlisting}[language=bash]
make run
\end{lstlisting}
构建并运行程序。

\begin{lstlisting}[language=bash]
make -n
\end{lstlisting}
只打印将要执行的命令，不真正执行。

\begin{lstlisting}[language=bash]
make -j4
\end{lstlisting}
使用 4 个并行任务加速构建。
\end{frame}

%------------------------------------------------
\begin{frame}[fragile,shrink]{调试技巧}
\textbf{1. 查看 Make 会执行什么}

\begin{lstlisting}[language=bash]
make -n
\end{lstlisting}

\textbf{2. 强制重新构建}

\begin{lstlisting}[language=bash]
make clean
make
\end{lstlisting}

\textbf{3. 显示变量值}

可以在 Makefile 中临时加入：

\begin{lstlisting}[language=make]
print:
	echo $(CC)
	echo $(CFLAGS)
	echo $(OBJS)
\end{lstlisting}

执行：

\begin{lstlisting}[language=bash]
make print
\end{lstlisting}
\end{frame}

%------------------------------------------------
\begin{frame}{常见错误}
\begin{block}{错误 1：missing separator}
原因通常是命令前面用了空格，而不是 Tab。

\[
\text{target: deps}
\]
\[
\text{Tab + command}
\]
\end{block}

\begin{block}{错误 2：No rule to make target}
Make 找不到生成某个依赖文件的规则。

可能原因：

\begin{itemize}
  \item 文件名写错
  \item 依赖文件不存在
  \item 规则没有定义
\end{itemize}
\end{block}
\end{frame}

%------------------------------------------------
\begin{frame}{Make 的工作流程}
Make 大致按照如下流程工作：

\begin{enumerate}
  \item 读取 \texttt{Makefile}
  \item 找到默认目标或指定目标
  \item 分析目标的依赖关系
  \item 判断依赖文件是否比目标文件更新
  \item 如果需要，执行对应命令
  \item 递归处理依赖目标
\end{enumerate}

\begin{block}{一句话总结}
Make 不是简单地顺序执行脚本，而是根据依赖关系决定是否执行命令。
\end{block}
\end{frame}

%------------------------------------------------
\section{课堂操作练习}

\begin{frame}[fragile]{练习：创建一个简单项目}
第一步：创建目录。

\begin{lstlisting}[language=bash]
mkdir make-demo
cd make-demo
\end{lstlisting}

第二步：创建三个源文件。

\begin{lstlisting}[language=bash]
touch main.c hello.c hello.h
\end{lstlisting}

第三步：创建 Makefile。

\begin{lstlisting}[language=bash]
touch Makefile
\end{lstlisting}

第四步：输入前面示例中的代码和 Makefile。
\end{frame}

%------------------------------------------------
\begin{frame}[fragile,shrink]{练习：执行常用操作}
构建项目：

\begin{lstlisting}[language=bash]
make
\end{lstlisting}

运行程序：

\begin{lstlisting}[language=bash]
make run
\end{lstlisting}

清理生成文件：

\begin{lstlisting}[language=bash]
make clean
\end{lstlisting}

观察增量构建：

\begin{lstlisting}[language=bash]
make
make
touch hello.c
make
\end{lstlisting}

思考：

\begin{itemize}
  \item 第二次 \texttt{make} 为什么没有重新编译？
  \item \texttt{touch hello.c} 后为什么会重新编译？
\end{itemize}
\end{frame}

%------------------------------------------------
\section{总结}

\begin{frame}[fragile]{总结}
本讲介绍了：

\begin{itemize}
  \item Make 的作用：自动化构建
  \item Makefile 的基本语法
  \item 目标、依赖和命令
  \item 变量的使用
  \item \texttt{.PHONY} 伪目标
  \item 自动变量：\texttt{\$@}、\texttt{\$<}、\texttt{\$\textasciicircum\{\}}
  \item 模式规则：\texttt{\%.o: \%.c}
  \item 常用命令：\texttt{make}、\texttt{make clean}、\texttt{make -n}、\texttt{make -j}
\end{itemize}

\begin{block}{核心记忆}
Makefile = 目标 + 依赖 + 命令
\end{block}
\end{frame}

%------------------------------------------------
\begin{frame}[fragile]{推荐后续学习}
可以继续学习：

\begin{itemize}
  \item 自动生成头文件依赖：\texttt{gcc -MMD -MP}
  \item 多目录项目 Makefile
  \item 静态库和动态库构建
  \item 与 CMake 的关系
  \item GNU Make 官方文档
\end{itemize}

\end{document}